% ===================================================================== % sections/intro_long.tex — long version (~3-4 pages) % =====================================================================

\section{Introduction} \label{sec:intro}

Indonesia entered the 1990s with one of Southeast Asia's highest maternal mortality ratios. In 1990 roughly 450 women died for every 100{,}000 live births \autocite{unicef}, a rate driven largely by unattended deliveries in rural communities where the only birth attendants were untrained traditional midwives (\emph{dukun bayi}). Mothers without access to skilled care faced the standard catastrophic complications of childbirth, obstructed labour, postpartum haemorrhage, eclampsia, that are largely preventable when a trained provider is present at the delivery \autocite{geefhuysen1999,makowiecka2008}. In 1989 the Indonesian government responded with a national program, \emph{Bidan di Desa} (``Village Midwife''), that trained and deployed tens of thousands of midwives, one per village, across rural Indonesia; by 1996 roughly \num{54000} village midwives had been placed, each serving the village on a guaranteed three-year government salary and permitted thereafter to continue in private practice in the same community \autocite{gani1996}. Because the services the midwives provided extended beyond delivery itself, antenatal monitoring, postnatal visits, family planning, nutrition counselling, and iron and vitamin A supplementation \autocite{frankenberg2005,makowiecka2008},the program is naturally interpreted as a bundle of early-life interventions targeted at the prenatal and early-postnatal windows that the developmental-medicine literature has identified as critical for long-run human capital formation \autocite{barker1995,duflo2001,maccini2009,majid2015}.

This paper asks whether early-life exposure to the program, in utero or during the first three years of life, left detectable effects on adult human-capital outcomes a generation later. The empirical setting makes the question answerable because the program's rollout was staggered across communities over nearly two decades, so the same calendar cohort of children received treatment at different ages depending on when a midwife first arrived in their village; and because the Indonesia Family Life Survey follows cohorts born in 1984--1999 into wave~5 (2014), when they are aged 15--30 and adult human-capital outcomes can be measured. Following the recent early-life-environments literature on socioemotional development \autocite{casepaxson2008,premand2022,sorrenti2025,webb2024,headstart}, the outcomes are collected into four inverse-covariance-weighted composite indices following \textcite{anderson2008}: physical health, cognition, depressive symptoms (CESD-10), and socioemotional traits (Big-Five Inventory).

The findings organise around three patterns. First, under the heterogeneity-robust Callaway--Sant'Anna estimator, which compares each treated cohort's outcome to a clean never-treated comparison group and therefore sidesteps the staggered-rollout biases that \textcite{goodmanbacon2021} documents for the canonical two-way-fixed-effects estimator, early-life exposure raises the adult health index by 0.31 standard deviations and the effect is statistically significant at the five-percent level. The point estimate is concentrated in the in-utero and age 0--3 event-time windows emphasised by the critical-periods literature. The same specification produces positive but statistically imprecise estimates on the cognition and depression indices; the socioemotional index is the one domain on which the effect is not positive. Second, the canonical two-way-fixed-effects estimator produces a near-zero point estimate on the same sample, a gap that is itself informative: it is consistent with \textcite{goodmanbacon2021}'s demonstration that already-treated units enter the effective control group for later-treated ones under the pooled estimator and can push the coefficient toward zero when treatment effects are heterogeneous across cohorts, and it illustrates why a paper on this program written without the heterogeneity-robust correction is liable to under-report the program's true human-capital effect. Third, a mechanism analysis of mother-reported pregnancy-history data traces a non-trivial share of the health-index effect to a higher probability of skilled birth attendance, which is the largest first-stage coefficient in the candidate mediator set and which is the component of the program's service bundle most directly delivered by the village midwife.

The pattern documented here departs from both the more cautious intermediate-horizon findings of \textcite{triyana2016}, who reported that most adult-health effects of the program did not persist ten years after launch, and the more optimistic long-horizon findings of \textcite{ahsan2021}, who report in-utero-exposure effects on adult health, cognition, and economic outcomes among men using a pooled specification that does not distinguish between the two-way-fixed-effects and heterogeneity-robust estimates. The interpretation is that once the 1984--1999 cohort window is combined with the correct staggered-rollout estimator, the program's effect on adult health is real and of a meaningful magnitude, but the literature's widely cited pooled-coefficient estimates have been contaminated by the staggered-rollout bias that this paper's Callaway--Sant'Anna column corrects. The contribution of the paper, against that backdrop, is threefold. The outcome set is extended to adult mental health and socioemotional traits, which the early-life-environments literature has identified as plausible long-run margins \autocite{webb2024} and which no prior study of this program has measured directly. The main results are reported under the heterogeneity-robust Callaway--Sant'Anna and de Chaisemartin--D'Haultf{\oe}uille estimators, with the pooled two-way-fixed-effects estimator retained only as a baseline for comparison with the earlier literature. The paper closes with a mechanism section that decomposes the reduced-form health effect into a first stage on three mother-reported mediators and a mediator-conditional effect on each adult outcome, which is the first such decomposition for this program using the pregnancy-history data that the IFLS itself provides.

The remainder of the paper is structured as follows. Section~\ref{sec:background} reviews the program's institutional setting. Section~\ref{sec:literature} positions the study against the early-life-shocks, program-specific, and socioemotional literatures. Section~\ref{sec:data} describes the data and sample construction. Section~\ref{sec:empirical} develops the identification strategy. Section~\ref{sec:results} presents the headline estimates and event-study evidence. Section~\ref{sec:robustness} reports the robustness suite, including the within-mother specification. Section~\ref{sec:heterogeneity} reports the gender and age-window heterogeneity. Section~\ref{sec:mechanisms} reports the mechanism analysis. Section~\ref{sec:discussion} interprets the findings and discusses limitations. Section~\ref{sec:conclusion} concludes.
